\begin{tikzpicture}
  \node[anchor=south west, inner sep=0] (polyimage) at (0,0)
    {\includegraphics[width=0.84\linewidth]{figures_zh/ch21/polysemanticity_graphics.pdf}};
  \begin{scope}[x={(polyimage.south east)},y={(polyimage.north west)}]
    \node[font=\small\sffamily\bfseries] at (0.526,0.967) {神经元激活空间中的多义性};
    \node[font=\footnotesize\sffamily] at (0.527,0.034) {神经元 A 的激活};
    \node[rotate=90, font=\footnotesize\sffamily] at (0.041,0.503) {神经元 B 的激活};
    \node[anchor=west, font=\footnotesize\sffamily\bfseries] at (0.838,0.275) {概念 1};
    \node[anchor=west, font=\footnotesize\sffamily\bfseries] at (0.259,0.699) {概念 2};
    \node[anchor=west, font=\footnotesize\sffamily\bfseries] at (0.722,0.738) {概念 3};
    \node[anchor=west, font=\footnotesize\sffamily\bfseries] at (0.335,0.468) {概念 4};
    \node[anchor=west, font=\footnotesize\sffamily\bfseries] at (0.645,0.391) {概念 5};
    \node[anchor=north west, inner sep=0, text width=0.420\linewidth,
          font=\footnotesize\sffamily, align=left] at (0.162,0.871) {
      \textbf{多义性：}\\
      单个神经元（坐标轴）参与表征多个不同概念。
      这些概念以激活空间中的“方向”分布，
      而不是与各个坐标轴完全对齐。
    };
  \end{scope}
\end{tikzpicture}
